\chapter{Experimental Setup}\label{chap:setup}

This chapter lists the configurations of all experiments reported in Chapter~\ref{chap:results}. The plaintext benchmark, the historical bounded secure runs and the two full-size secure runs (Secure BP and Secure FF) are separate experimental regimes and are described in separate tables. Measurement definitions are given in Section~\ref{sec:method-measure}.

\section{Data}\label{sec:setup-data}
All experiments use MNIST (28$\times$28 grey-scale digits, 784 input values, 10 classes). Table~\ref{tab:setup-data} lists the splits.

\begin{table}[ht]
\centering
\small
\begin{tabular}{@{}>{\raggedright\arraybackslash}p{3.6cm}>{\raggedright\arraybackslash}p{4.4cm}>{\raggedright\arraybackslash}p{5.8cm}@{}}
\toprule
Regime & Train / validation / test & Source of the examples \\
\midrule
Plaintext benchmark (Plain BP, Plain FF) & 5{,}000 / 1{,}000 / 1{,}000 & Train and validation: first 5{,}000 and next 1{,}000 positions of a seed-42 permutation of the official training set; test: positions 1{,}000--1{,}999 of a seed-42 permutation of the official test set \\
Full-size Secure BP & 5{,}000 / 1{,}000 / 1{,}000 & Same split function and indices as the plaintext benchmark \\
Full-size Secure FF & 5{,}000 / 1{,}000 / 1{,}000 & Same split function and indices as the plaintext benchmark (checked by an independent verifier against the dataset and the Secure BP bundle); seed 42 \\
Bounded Secure BP (A0, A1, BP-COMP-001) & 90 / 30 / 30 & Class-balanced subset, selection seed 42; 30 / 10 / 10 per party \\
\bottomrule
\end{tabular}
\caption{Data splits of the three experimental regimes.}
\label{tab:setup-data}
\end{table}

\section{Plaintext Benchmark}\label{sec:setup-plain}
Table~\ref{tab:setup-plain} gives the configuration of the plaintext benchmark (identifier PT-BENCH-v1, profile \texttt{project\_baseline\_standardized}).

\begin{table}[ht]
\centering
\small
\begin{tabular}{@{}>{\raggedright\arraybackslash}p{3.6cm}>{\raggedright\arraybackslash}p{5.0cm}>{\raggedright\arraybackslash}p{5.2cm}@{}}
\toprule
Item & Plain BP & Plain FF \\
\midrule
Architecture & 784--128--128--10, ReLU, linear output & 784--128--128, ReLU, L2 normalisation between layers \\
Parameters & 118{,}282 & 116{,}992 \\
Initialisation & PyTorch default, \texttt{manual\_seed(seed)} & PyTorch default, \texttt{manual\_seed(seed)} \\
Objective & Softmax cross-entropy & Per-layer softplus goodness loss, threshold 2.0 \\
Label handling & Target class & One-hot label in first 10 pixels; negatives with a uniformly random wrong label \\
Optimizer & \ac{sgd}, learning rate 0.05 & Per-layer \ac{sgd}, learning rate 0.01 \\
Batch size; epochs & 128; 5 & 128; 5 per layer, layers trained in sequence \\
Optimizer updates & 200 & 400 \\
Example presentations & 25{,}000 & 100{,}000 (positive and negative) \\
Inference & One forward pass, argmax of logits & Ten candidate-label passes, highest goodness summed over both layers \\
Checkpoint & Best validation accuracy over epochs & Best validation accuracy over layer-2 epochs \\
Seeds; repetitions & 42--46; 3 timing repetitions per seed & 42--46; 3 timing repetitions per seed \\
\bottomrule
\end{tabular}
\caption{Configuration of the plaintext benchmark PT-BENCH-v1.}
\label{tab:setup-plain}
\end{table}

\section{Historical Bounded Secure BP Runs}\label{sec:setup-bounded}
Table~\ref{tab:setup-secure} lists the configuration of the bounded secure \ac{bp} runs.

\begin{table}[ht]
\centering
\small
\begin{tabular}{@{}>{\raggedright\arraybackslash}p{4.3cm}>{\raggedright\arraybackslash}p{9.6cm}@{}}
\toprule
Item & Setting \\
\midrule
Framework, protocol & MP-SPDZ v0.4.4; Rep3 Ring; 3 parties; semi-honest, honest majority \\
Arithmetic & 64-bit ring; \texttt{sfix(16,31)}; probabilistic truncation and \texttt{use\_split(3)} for A1 and BP-COMP-001 (the A0 log suggests that these options were not enabled for A0) \\
Program & Scalar loops (\texttt{bp\_train.mpc}) \\
Model & 784--8--10 \ac{mlp}, ReLU hidden layer, linear output layer, 6{,}370 parameters \\
Data & MNIST subset: 90 train / 30 validation / 30 test; 30/10/10 per party \\
Training & Batch size 3 (one per party); 30 batches per epoch; 3 epochs; 90 updates; learning rate 0.05; no momentum or weight decay \\
Loss, output & Stable softmax cross-entropy; aggregate-only outputs \\
Historical runs (A0, A1) & Class-major batches; deterministic initialisation $(i-j)/1000$, zero biases \\
BP-COMP-001 & Per-party shuffled order (ordering seed 42042); He (hidden) and Xavier uniform (output) initialisation (seed 42); zero biases \\
Metrics & Validation/test correct out of 30; losses; MP-SPDZ time; global communication; compiler VM rounds \\
Repetitions & One run per configuration \\
Execution host & All three parties on one host (localhost) \\
\bottomrule
\end{tabular}
\caption{Configuration of the bounded secure \ac{bp} experiments.}
\label{tab:setup-secure}
\end{table}

\section{Full-Size Secure BP Run}\label{sec:setup-full}
Table~\ref{tab:setup-full} lists the configuration of the full-size run. All items were checked against the program source, the configuration file \texttt{configs/\allowbreak bp\_\allowbreak matched.yaml}, the input-bundle metadata and the run log.

\begin{table}[ht]
\centering
\small
\begin{tabular}{@{}>{\raggedright\arraybackslash}p{4.3cm}>{\raggedright\arraybackslash}p{9.6cm}@{}}
\toprule
Item & Setting \\
\midrule
Experiment & SECURE-BP-MATCHED-v1, batched-matrix program, seed 42 \\
Program & \texttt{mpc/programs/bp/bp\_matched\_matrix.mpc} \\
Framework, protocol & MP-SPDZ v0.4.4 (commit \texttt{113c4e47}); Rep3 Ring; 3 parties; semi-honest, honest majority \\
Arithmetic & 64-bit ring; \texttt{sfix(16,31)}; probabilistic truncation; \texttt{use\_split(3)} \\
Compile and run command & \texttt{compile-run.py -l -M -E ring bp\_matched\_matrix -- -IF \textless{}bundle\textgreater{}/Input} \\
Model & 784--128--128--10, ReLU hidden layers, linear output layer, 118{,}282 parameters \\
Initialisation & PyTorch default after seed 42, transposed to input-by-output layout; input by party 0 \\
Data & 5{,}000 / 1{,}000 / 1{,}000 (Table~\ref{tab:setup-data}); example owner = position modulo 3 \\
Order & Baseline seed-42 data loader order, state continued across epochs \\
Training & Batch size 128; 39 batches of 128 and one of 8 per epoch (40 batches); 5 epochs; 200 updates; \ac{sgd}, learning rate 0.05, no momentum or weight decay; gradients averaged over the actual batch size \\
Loss & Stable softmax cross-entropy \\
Evaluation & Final (epoch-5) model; argmax of logits; chunks of 100 \\
Disclosure & Mean training loss per epoch; validation and test mean loss and correct count \\
Repetitions & One run (seed 42); no timing repetitions \\
Date & 15 November 2026 (run directory timestamp 23:49:46 local time) \\
\bottomrule
\end{tabular}
\caption{Configuration of the full-size secure \ac{bp} run.}
\label{tab:setup-full}
\end{table}

\section{Full-Size Secure FF Run}\label{sec:setup-sff}
Table~\ref{tab:setup-sff} lists the configuration of the full-size Secure \ac{ff} run. The items were taken from the configuration file \texttt{configs/\allowbreak ff\_\allowbreak secure\_\allowbreak matched.yaml}, the program header and source, the command record and the run output. Where the Secure \ac{bp} and Secure \ac{ff} settings differ, the difference is intended (learning rule) or is a documented confound (Section~\ref{sec:method-sff}).

\begin{table}[htbp]
\centering
\small
\begin{tabular}{@{}>{\raggedright\arraybackslash}p{4.3cm}>{\raggedright\arraybackslash}p{9.6cm}@{}}
\toprule
Item & Setting \\
\midrule
Experiment & SECURE-FF-MATCHED-v1, profile \texttt{full}, seed 42; run directory \texttt{seed42\_run\_20261122\_212319} \\
Program & \texttt{mpc/programs/ff/ff\_matched.mpc} \\
Framework, protocol & MP-SPDZ v0.4.4 (commit \texttt{113c4e47}); Rep3 Ring; 3 parties; semi-honest, honest majority \\
Arithmetic & 64-bit ring; \texttt{sfix(16,31)}; default probabilistic truncation \\
Compile command & \texttt{compile.py -M -E ring ff\_matched} (no \texttt{-l}) \\
Run command & \texttt{Scripts/ring.sh ff\_matched -IF \textless{}bundle\textgreater{}/Input}, three party processes on \texttt{localhost} \\
Model & 784--128--128, ReLU, 116{,}992 parameters; no output layer \\
Initialisation & Plain FF model after \texttt{random.seed} and \texttt{torch.manual\_seed} with seed 42, transposed to input-by-output layout; input by party 0 \\
Data & 5{,}000 / 1{,}000 / 1{,}000 (Table~\ref{tab:setup-data}); record owner = position modulo 3 \\
Order & One seeded data loader iterated ten times; state continued across layers \\
Local objective & $\frac1B\sum\mathrm{softplus}(\theta-g^{+})+\frac1B\sum\mathrm{softplus}(g^{-}-\theta)$; $\theta=2.0$; goodness = sum of squared ReLU activations before normalisation \\
Normalisation & Stable max-scaled L2 normalisation between the layers \\
Negative labels & Uniformly random incorrect label, generated in \ac{mpc} (ten rejection-sampling tries of eight secret bits) \\
Optimizer & Per-layer \ac{sgd}, learning rate 0.01, no momentum or weight decay; update $\eta/B$ applied as $a\cdot2^{-s}$ \\
Batches & Batch size 128; 39 batches of 128 and one of 8 per epoch (40 batches); every example of a batch is presented once with its correct label (positive) and once with an incorrect label (negative), in one stacked matrix \\
Epochs and updates & 5 epochs for layer 1, then 5 epochs for layer 2 (layer 1 frozen); 200 updates per layer, 400 in total \\
Inference & Ten candidate labels, chunks of 50 samples, score = goodness of layer 1 + goodness of layer 2, argmax with lowest index on ties \\
Checkpoint & Layer-2 epoch with the strictly highest validation correct count (first epoch always selected), restored before the single test evaluation \\
Disclosure & Weighted loss per layer-epoch; validation correct count per layer-2 epoch; selected epoch; test correct count; sampler failure count \\
Repetitions & One run (seed 42); no timing repetitions \\
Date & 22 November 2026 (command record: 20:23:19 UTC) \\
\bottomrule
\end{tabular}
\caption{Configuration of the full-size Secure \ac{ff} run. The rehearsal run used the \texttt{mini} profile (20/10/10 examples, batch 8, two epochs per layer) and is not an experiment.}
\label{tab:setup-sff}
\end{table}

Table~\ref{tab:setup-diff} lists the configuration differences between the two secure programs that matter for the interpretation of the preliminary comparison.

\begin{table}[htbp]
\centering
\small
\begin{tabular}{@{}>{\raggedright\arraybackslash}p{3.3cm}>{\raggedright\arraybackslash}p{5.2cm}>{\raggedright\arraybackslash}p{5.2cm}@{}}
\toprule
Item & Secure BP & Secure FF \\
\midrule
Architecture; parameters & 784--128--128--10; 118{,}282 & 784--128--128; 116{,}992 \\
Learning rate & 0.05 & 0.01 \\
Updates; presentations & 200; about 25{,}000 & 400; about 100{,}000 (positive and negative) \\
Evaluation & Final epoch-5 model; one forward pass per sample & Validation-selected layer-2 epoch; ten candidate passes per sample \\
Compiler invocation & \texttt{compile-run.py -l -M -E ring} & \texttt{compile.py -M -E ring} \\
External wall time covers & Compilation and execution & Execution only; compilation timed separately \\
Peak memory covers & Largest process of the timed tree (compilation included) & Largest of three party processes (compilation excluded) \\
\bottomrule
\end{tabular}
\caption{Differences between the Secure BP and Secure FF configurations.}
\label{tab:setup-diff}
\end{table}

\section{Execution Environment}\label{sec:setup-env}
Table~\ref{tab:setup-env} records what is known about the execution environment. Only information present in the run artifacts is treated as recorded. The host description of the secure runs was obtained by inspecting the machine that holds the run directories on 16 November 2026; it is not part of the run log, so it is reported as the likely, but unverified, environment of the run. The Secure FF command record shows the same arrangement (three party processes on \texttt{localhost}, \texttt{-h localhost}) and does not contain a hardware description; the host of the Secure FF run is therefore likewise not verified from its artifacts.

\begin{table}[ht]
\centering
\small
\begin{tabular}{@{}>{\raggedright\arraybackslash}p{3.6cm}>{\raggedright\arraybackslash}p{4.8cm}>{\raggedright\arraybackslash}p{5.4cm}@{}}
\toprule
Item & Plaintext benchmark & Secure runs \\
\midrule
Software & Python 3.14.4; PyTorch 2.14.1+cpu; 4 threads (recorded) & MP-SPDZ v0.4.4, commit \texttt{113c4e47} (recorded) \\
Party placement & Not applicable & Three party processes on one host, \texttt{-h localhost} (recorded) \\
Network & Not applicable & Local loopback connections (recorded command line); no network emulation recorded; no LAN or WAN measurement \\
Processor & Not recorded & Not recorded in the log; host inspected after the run: Intel Core i5-11320H, 8 logical CPUs visible \\
Memory & Not recorded & Not recorded in the log; host inspected after the run: about 11\,GiB visible \\
Operating system & Not recorded & Not recorded in the log; host inspected after the run: Ubuntu 26.04.1 LTS under WSL2 (kernel 6.6.87.2) \\
\bottomrule
\end{tabular}
\caption{Execution environment. ``Recorded'' means present in the experiment artifacts.}
\label{tab:setup-env}
\end{table}

Because hardware and network were not recorded in the run artifacts and all parties ran on one host, the runtime and communication figures describe these runs only and are not transferable to other environments.
